% BEGIN VOL07-EXPANSION VII12
\section{Supplementary worked examples}
\label{sec:vii12-pedagogy}

\begin{example}[Arc-length reparametrization of a parabola]\label{ex:vii12-ped-01}
For \(\gamma(t)=(t,t^2)\), the speed is \(\sqrt{1+4t^2}\).  The arc-length function
\[
s(t)=\int_0^t\sqrt{1+4u^2}\,du
\]
has strictly positive derivative, so it is locally invertible.  Reparametrizing by \(s\) produces a unit-speed curve and separates geometric length from the arbitrary choice of parameter.
\end{example}

\begin{example}[A helix has constant speed]\label{ex:vii12-ped-02}
The curve \(\gamma(t)=(a\cos t,a\sin t,bt)\) has
\[
\|\gamma'(t)\|=\sqrt{a^2+b^2}.
\]
Thus \(s=\sqrt{a^2+b^2}\,t\) is an arc-length parameter up to an additive constant.  Helices provide a clean model in which tangent direction changes while speed remains constant.
\end{example}

\begin{example}[Regularity fails at a cusp]\label{ex:vii12-ped-03}
The parametrized curve \(\gamma(t)=(t^2,t^3)\) traces a cusp.  Its derivative vanishes at \(t=0\), so it is not regular there.  The image may still be a perfectly meaningful subset, but the parametrization does not provide a one-dimensional smooth immersion at the cusp.
\end{example}

\section{Graded supplementary exercises}
The following exercises add a deliberate progression from concrete calculation to structural proof, hypothesis testing, application, and synthesis.

\subsection*{Standard computations and constructions}

\begin{exercise}[Speed of a helix]\label{exr:vii12-ped-01}
Compute the speed of \(\gamma(t)=(2\cos t,2\sin t,3t)\).
\end{exercise}
\begin{hint}
Differentiate and take the Euclidean norm.
\end{hint}
\begin{solution}
\(\gamma'=(-2\sin t,2\cos t,3)\), so the speed is \(\sqrt{4+9}=\sqrt{13}\).
\end{solution}

\begin{exercise}[Length of a circle arc]\label{exr:vii12-ped-02}
Find the length of \(\gamma(t)=(R\cos t,R\sin t)\), \(0\le t\le\theta\).
\end{exercise}
\begin{hint}
The speed is constant.
\end{hint}
\begin{solution}
\(\|\gamma'\|=R\), hence the length is \(R\theta\) for \(\theta\ge0\).
\end{solution}

\begin{exercise}[Regularity test]\label{exr:vii12-ped-03}
Determine where \(\gamma(t)=(t^2,t^3)\) is regular.
\end{exercise}
\begin{hint}
Find where \(\gamma'(t)\neq0\).
\end{hint}
\begin{solution}
\(\gamma'=(2t,3t^2)\), which vanishes only at \(t=0\); the curve is regular elsewhere.
\end{solution}

\begin{exercise}[Unit-speed line]\label{exr:vii12-ped-04}
Give a unit-speed parametrization of the line through \(p\) in direction \(v\neq0\).
\end{exercise}
\begin{hint}
Normalize the direction vector.
\end{hint}
\begin{solution}
\(\gamma(s)=p+s\,v/\|v\|\) is unit speed.
\end{solution}

\begin{exercise}[Reparametrized speed]\label{exr:vii12-ped-05}
If \(\tilde\gamma(u)=\gamma(\phi(u))\), express \(\|\tilde\gamma'(u)\|\).
\end{exercise}
\begin{hint}
Use the chain rule.
\end{hint}
\begin{solution}
\(\|\tilde\gamma'(u)\|=\|\gamma'(\phi(u))\|\,|\phi'(u)|\).
\end{solution}

\subsection*{Proofs}

\begin{exercise}[Length invariance]\label{exr:vii12-ped-06}
Prove that curve length is invariant under orientation-preserving \(C^1\) reparametrization.
\end{exercise}
\begin{hint}
Use the substitution \(t=\phi(u)\).
\end{hint}
\begin{solution}
The speed picks up the factor \(\phi'(u)>0\), which is exactly canceled by the change of variables in the integral.
\end{solution}

\begin{exercise}[Local arc-length parameter]\label{exr:vii12-ped-07}
Prove that every regular curve admits a local arc-length parameter.
\end{exercise}
\begin{hint}
Use \(s'(t)=\|\gamma'(t)\|>0\).
\end{hint}
\begin{solution}
The arc-length function has nonzero derivative, so the inverse function theorem provides a local inverse and the reparametrized curve has unit speed.
\end{solution}

\begin{exercise}[Unit tangent orthogonality]\label{exr:vii12-ped-08}
For a unit-speed curve, prove \(T\cdot T'=0\).
\end{exercise}
\begin{hint}
Differentiate \(T\cdot T=1\).
\end{hint}
\begin{solution}
Differentiation gives \(2T\cdot T'=0\).
\end{solution}

\begin{exercise}[Rigid motions preserve regularity and length]\label{exr:vii12-ped-09}
Prove Euclidean rigid motions preserve regular curves and their lengths.
\end{exercise}
\begin{hint}
Differentiate \(A\gamma+b\) with \(A\) orthogonal.
\end{hint}
\begin{solution}
The derivative becomes \(A\gamma'\), whose norm equals \(\|\gamma'\|\); nonzero derivatives and length integrals are preserved.
\end{solution}

\subsection*{Counterexamples and hypothesis tests}

\begin{exercise}[Injectivity is part of regularity]\label{exr:vii12-ped-10}
Test the claim with \(\gamma(t)=(\cos t,\sin t)\) on \(\mathbb R\).
\end{exercise}
\begin{hint}
Check the derivative and periodicity separately.
\end{hint}
\begin{solution}
False.  The derivative never vanishes, so the curve is regular although the parametrization is not injective.
\end{solution}

\begin{exercise}[Every parametrization of a smooth curve is regular]\label{exr:vii12-ped-11}
Test by reparametrizing a line with \(t\mapsto t^3\).
\end{exercise}
\begin{hint}
Differentiate at zero.
\end{hint}
\begin{solution}
False.  \(\gamma(t)=(t^3,0)\) has zero derivative at \(0\) despite tracing a straight line.
\end{solution}

\begin{exercise}[Arc length depends on speed choice]\label{exr:vii12-ped-12}
Test whether two regular reparametrizations of the same oriented curve have different lengths.
\end{exercise}
\begin{hint}
Use change of variables.
\end{hint}
\begin{solution}
False.  Length is invariant under regular reparametrization even though pointwise speed changes.
\end{solution}

\subsection*{Applications and investigations}

\begin{exercise}[Motion along a path]\label{exr:vii12-ped-13}
Interpret the speed and length of a curve as kinematic quantities.
\end{exercise}
\begin{hint}
Think of \(t\) as time.
\end{hint}
\begin{solution}
\(\|\gamma'(t)\|\) is instantaneous speed, and its time integral is total distance traveled along the parametrized path.
\end{solution}

\begin{exercise}[Sampling a curve uniformly]\label{exr:vii12-ped-14}
Why is arc-length parametrization useful for numerical sampling?
\end{exercise}
\begin{hint}
Equal parameter increments should correspond to equal geometric distances to first order.
\end{hint}
\begin{solution}
For unit speed, small equal increments in \(s\) have equal first-order arc length, avoiding clustering caused solely by parameter distortion.
\end{solution}

\subsection*{Challenge problems}

\begin{exercise}[Closed unit-speed curve]\label{exr:vii12-ped-15}
If a unit-speed curve has period \(L\), show the length of one period is \(L\).
\end{exercise}
\begin{hint}
Integrate the constant speed \(1\).
\end{hint}
\begin{solution}
The length is \(\int_0^L1\,ds=L\).
\end{solution}

\begin{exercise}[Regularity under reparametrization]\label{exr:vii12-ped-16}
Give the precise condition on a change of parameter \(\phi\) that preserves regularity.
\end{exercise}
\begin{hint}
Use \((\gamma\circ\phi)'=\gamma'(\phi)\phi'\).
\end{hint}
\begin{solution}
For a regular \(\gamma\), the reparametrized curve is regular wherever \(\phi'\neq0\); a regular reparametrization requires a smooth local diffeomorphism of parameter intervals.
\end{solution}

% END VOL07-EXPANSION VII12
